\begin{figure}[!t]
\centering
\begin{subfigure}[t]{0.322\linewidth}
\centering\includegraphics[width=\linewidth]{figures/pdf/fig3a_context_trajectory.pdf}
\caption{Canonical context trajectories.}
\end{subfigure}\hfill
\begin{subfigure}[t]{0.322\linewidth}
\centering\includegraphics[width=\linewidth]{figures/pdf/fig3b_deletion_trajectory.pdf}
\caption{Canonical deletion trajectories.}
\end{subfigure}\hfill
\begin{subfigure}[t]{0.332\linewidth}
\centering\includegraphics[width=\linewidth]{figures/pdf/fig3c_balanced_effects.pdf}
\caption{Balanced endpoint changes.}
\end{subfigure}
\caption{\textbf{Controlled stress does not uniformly increase the coverage gap.} Curves use fixed canonical-interface fits across five doses. Endpoint changes use separately fitted balanced-interface endpoint designs, with pointwise 95\% common-question bootstrap intervals. Positive changes indicate greater reconstruction error. D denotes deletion; I denotes irrelevant context.}
\label{fig:stress}
\end{figure}

\section{Coverage under stress and response balancing}
\label{sec:stress}

The identity of the peers is fixed, but the response relation can change with the input. Irrelevant-context injection prepends unrelated material; content deletion removes selected words from the question stem while preserving the options. A single combination is fitted across the declared conditions of each family and held fixed during evaluation. Changes in its residual reveal where that combination ceases to track the target.

\paragraph{Opposite directions under matched stress.}
At the maximum irrelevant-context dose, the balanced reference-answer gap increases by $0.01435$ for Qwen2.5 and $0.03250$ for phi-4. It decreases by $0.00894$ for Mistral and $0.01026$ for OLMo. All four directions hold in every split and have common-question bootstrap intervals excluding zero (Figure~\ref{fig:stress}c). Under deletion, Mistral and OLMo also decrease, by $0.02183$ and $0.01565$. The same perturbation therefore makes one target harder to reconstruct and another easier.

The canonical-order curves show the full dose dependence (Figure~\ref{fig:stress}a--b). Balanced endpoint fits then examine the high-dose changes under a second presentation protocol. The curves and endpoint comparisons retain their own response definitions and fitting designs. The complete-option diagnostic also gives opposite irrelevant-context directions for Qwen2.5 and Mistral (Appendix~\ref{app:vector}).



\paragraph{Presentation changes magnitudes more often than directions.}
The mean rotation-to-average TV diagnostic is $0.338$, indicating substantial variation in semantic option probabilities across presentations. Nevertheless, canonical and balanced endpoint analyses agree on 15 of 16 effect signs. Their eight-target rankings of family-averaged reconstruction errors have Spearman correlation $0.952$; the correlation across the 16 endpoint changes is $0.932$. These are complementary summaries of the error level and its response to stress.

Granite illustrates why both the effect and its uncertainty matter. Its canonical contractions remain negative after balancing, but the balanced intervals include zero. The resulting evidence supports the Qwen2.5, phi-4, Mistral, and OLMo irrelevant-context directions while distinguishing them from the weaker Granite effects. Thus presentation balancing refines which coverage changes are supported, rather than leaving every conclusion unchanged. Appendix~\ref{app:robustness} gives all endpoint estimates and the separate direct-generation diagnostics.
